\chapter{第二示例章}
\label{ch:02}

本章用于演示章文件的最小结构：一个章标题、若干节、正文、脚注与图。实际项目中，每章一个文件，文件名与原文章号对应。

\section{章节文件约定}

章节文件只包含正文，不写导言区，不写 \verb|\begin{document}|。标题用 \verb|\chapter|，节用 \verb|\section|；不要手写“第一章”，编号由模板自动生成。

\section{与原文对齐}

段落与原文保持一一对应。若中文需要合并或拆分段落，应在 \texttt{notes/decisions.md} 中记录位置与理由，便于审校时对照。

跨章引用用标签，不用固定数字：

\begin{quote}
详见第 \ref{ch:01} 章关于术语一致性的讨论。
\end{quote}

这样增删章节后，编号与交叉引用会自动更新。
